
\section{Move Semantics}

\begin{frame}[fragile]{Value Categories}
  Every \emph{expression} has a \emph{value category} (lvalue or rvalue).
  \begin{minted}{cpp}
    struct lorem { struct ipsum { int x; } y; int z; } l;
    int* ptr = &l.z;
    int arr[3] = {1, 2, 3};
  \end{minted}

  \begin{columns}[T]
    \column{0.48\textwidth}
    \textbf{lvalues} — has identity (an address)
    \begin{minted}{cpp}
l             // named variable
l.z           // member access
l.y.x         // nested member
*ptr          // dereference
arr[0]        // subscript
    \end{minted}
    Can appear on the \emph{left} side of assignment.\\
    Can take the address with \code{\&}.

    \column{0.48\textwidth}
    \textbf{rvalues} — no persistent identity
    \begin{minted}{cpp}
42            // literal
l.z + 1       // arithmetic result
intbox(3)     // temporary object
std::move(l)  // cast to rvalue
int{}         // value-initialized tmp
    \end{minted}
    Cannot appear on the left side of assignment.\\
    Cannot take the address with \code{\&}.
  \end{columns}
\end{frame}

\begin{frame}[fragile]{rvalue references}
  \begin{sqt:definition}
    An \emph{rvalue reference} is a reference that can bind to an \emph{rvalue}.
  \end{sqt:definition}
  \begin{sqt:result}
    When a ref is bound to a tmp, lifetime extended until the ref goes out of scope.\\
    i.e. the compiler postpones the dtor to the end of the scope.
  \end{sqt:result}

  \begin{minted}{cpp}
    /* 1.03 */
    {
        // tell compiler: we want to extend the tmp's lifetime until end of scope
        intbox&& rref = intbox(3);
        // play around with the temporary via the rvalue ref
        std::cout << rref.get() << std::endl;
        rref.get() *= 10;
        std::cout << rref.get() << std::endl;
    }
  \end{minted}
\end{frame}

\begin{frame}[fragile]{Arguments by rvalue ref}
  We already know that a function can take arguments by (lvalue) reference.\\
  It is also possible for a function to take arguments by rvalue reference.

  \begin{minted}{cpp}
/* 1.03 */
void foo(intbox& obj) { obj.data_ += 42; }
void foo(intbox&& obj) { obj.data_ = nullptr; }

int main() {
    // If you pass in an rvalue (e.g. the tmp result of the ctor function intbox(3)),
    // the rvalue ref extends the lifetime of the tmp until the function returns.
    foo(intbox(3));
}
  \end{minted}
\end{frame}

\begin{frame}[fragile]{Overload resolution}
  Given that an argument has two overloads, one taking lvalues and one taking rvalues,
  how does the compiler determine which overload to dispatch for each call-site?

  \begin{minted}{cpp}
    {
        intbox x{3};
        intbox&& rref = std::move(x);

        foo(x); // calls lvalue overload
        foo(rref); // also calls lvalue overload (weird)
        foo(intbox(3)); // calls rvalue overload
        foo(std::move(x)); // calls rvalue overload
        foo(std::move(rref)); // calls rvalue overload
    }
  \end{minted}

  \begin{sqt:result}
    The compiler picks the overload depending on the \emph{value category} of the argument.
  \end{sqt:result}
\end{frame}

\begin{frame}[fragile]{Solution: Move Construction}
  \begin{sqt:definition}
    A move constructor is a constructor that takes an rvalue reference
    to an existing object of the same type.
  \end{sqt:definition}

  \begin{minted}{cpp}
    struct intbox {
        /* ... */
        intbox(intbox&& other) : data_{std::exchange(other.data_, nullptr)} {
            std::cout << "move constructor" << std::endl;
        }
    };

    int main() {
        intbox a(42);
        intbox b = std::move(a); // b.data_ = a.data_; a.data_ = nullptr; (cheap operation)
        // a now points to nothing
        // ~b; // b's dtor frees the heap-allocated int
        // ~a; // a's destructor frees nullptr (no-op).
    }
  \end{minted}
\end{frame}

\begin{frame}[fragile]{Returning a local variable by moving}
  \begin{minted}{cpp}
    intbox product(intbox& a, intbox& b) {
        auto result = *(a.data_) * *(b.data_);
        intbox retval{result};
        return std::move(retval);
    }

    int main() {
        intbox a(5); intbox b(6);
        intbox c = product(a, b);
        // retval is move constructed into the return slot c (cheap operation)
        // retval is then destroyed (but it owns nothing at this point)
    }
  \end{minted}

  \begin{sqt:tip}
    RVO/NRVO (constructing directly in return slot)
    is preferable over constructing inside a local variable,
    and then \code{std::move} into the return slot.
  \end{sqt:tip}
\end{frame}

\begin{frame}[fragile]{Issue: Ownership}
  Suppose we have a class that handles a Unique instance of an object. 

  \begin{minted}{cpp}
    int main() {
      int* x = new int(4);

      uniqueIntbox a(x); // We do not want anyone else to be able to copy

      uniqueIntbox b(a); // What should we do here?
    }
  \end{minted}

  Here the raw int pointer is being binded to the uniqueIntbox class. This is a common
  pattern we see in C++ (examples include \mintinline{cpp}{std::unique_ptr}, and \mintinline{cpp}{std::unique_lock}). 
  
  What constructor should we allow \code{b} to be constructed with? If we allow for deep copies,
  then we violate the ``uniqueness" of the data held by \code{x}. If we allow shallow copies, then
  we get the same double free issue we had before. 
\end{frame}

\begin{frame}[fragile]{Solution: Move Semantics}
  In situations like this, we typically want to only allow the programmer to transfer
  ownership to the underlying data.

  \begin{minted}{cpp}
    struct uniqueIntbox {
      int *data_;

      intbox (int* x) : data_(x) {};
      ~intbox() { delete data_; }
      intbox (const intbox& other) = delete;
      intbox (intbox&& other) {
        this->data = other.data;
        other.data = nullptr;
      }
    }
  \end{minted}

  Here we disable to copy ctor entirely, and only allow objects to be constructed via r-value
  references to another \code{uniqueIntbox}.
\end{frame}

\begin{frame}[fragile]{Transferring Ownership}
  \begin{minted}{cpp}
    int main() {
      int* x = new int(4);
      
      uniqueIntbox a(x);
      uniqueIntbox b = std::move(a); // Trasnfer ownership of a to b
      uniqueIntbox c(b); // Compile Error
    }
  \end{minted}

  We cast \code{a} to an r-value before passing it into \code{b}'s constructor. Importantly,
  the l-value \code{a} is left in an invalid state after the move not undefined. Trying to 
  invoke the \code{c}'s copy constructor will prevent the program from compiling.
\end{frame}

\begin{frame}[fragile]{Transferring Ownership Trap}
  \code{std::move} will only attempt an r-value cast, but does not guarantee it.

  \begin{minted}{cpp}
    int main() {
      int* x = new int(4);

      const uniqueIntbox a(x);
      uniqueIntbox b = std::move(a); // Will attempt to copy
    }
  \end{minted}

  Here, since \code{x} is qualified is \code{const}, the r-value constructor cannot be used,
  and so the above code will not compile if the copy constructor is deleted. If there is a copy constructor
  available, the above code will silently invoke the copy constructor, which is likely not what was desired.

\end{frame}

\begin{frame}[fragile]{Ownership Trap II}
  What do you think will happen below:
  \begin{minted}{cpp}
    int main() {
      int* x = new int(4);
    
      uniqueIntbox a(x);
      uniqueIntbox b(x);
    }
  \end{minted}

  Who ``owns" the pointer? What will happen when you run this code?
\end{frame}


% % Used for
% \begin{frame}[fragile]{Non-Copyable Stuff}
%   Let's create a non-copyable int box


%   Lets add a move constructor to our wrapper class and delete the copy constructor
%   \begin{minted}{cpp}
%     Class gadgetWrapper {
%     public:
%       gadgetWrapper(gadget* rawGadget) : gadget_(rawGadget) {}

%       gadgetWrapper(const gadgetWrapper& other) = delete;

%       gadgetWrapper(gadgetWrapper&& other) {
%         this->gadget_ = other.gadget_;
%         other.gadget_ = nullptr;
%       }
      
%       ~gadgetWrapper() {
%         if (gadget_) delete gadget_;
%       }

%     private:
%       int* gadget_;
%     }
%   \end{minted}
% \end{frame}

% \begin{frame}[fragile]{Moving Ownership cont}
%   Now the previous example won't even compile!
%   \begin{minted}{cpp}
%   int main() {
%     gadget* c = new gadget(42);

%     gadgetWrapper d(c); // Binding c to d
%     gadgetWrapper e = d; // ERROR: Can't copy construct and steal data!
%   }
% \end{minted}
% \end{frame}

% \begin{frame}[fragile]{Solution: Move Semantics}
%   We ask the compiler to cast d into a temporary R-Value!
%   \begin{minted}{cpp}
%   int main() {
%     gadget* c = new gadget(42);

%     gadgetWrapper d(c); // Binding c to d
%     gadgetWrapper e = std::move(d); // Casts d into gadgetWrapper&& type

%     // d is in a valid state but owns nothing.

%     // ~e frees c, ~d does nothing avoiding the double free.
%   }

%   \end{minted}
% \end{frame}